% !TeX program = xelatex
% ============================================================
% 通用中文学术论文模板（biblatex + biber + GB/T 7714-2015）
% 编译：xelatex main.tex → biber main → xelatex main.tex ×2
% 用法：连同 refs-sample.bib 一起复制到 PROJECT_ROOT，改名为 refs.bib
% ============================================================
\documentclass[UTF8,a4paper,zihao=-4]{ctexart}

% ---------- 版面 ----------
\usepackage{geometry}
\geometry{a4paper,top=2.5cm,bottom=2.5cm,left=2.8cm,right=2.8cm}

% ---------- 数学与图表 ----------
\usepackage{amsmath,amssymb,bm}
\usepackage{graphicx}
\usepackage{booktabs}
\usepackage{caption}
\usepackage{enumitem}

% ---------- 颜色与链接 ----------
\usepackage{xcolor}
\definecolor{accent}{HTML}{1F4E79}

% ---------- 参考文献（默认：biblatex + biber + GB/T 7714） ----------
\usepackage[backend=biber,style=gb7714-2015]{biblatex}
\addbibresource{refs.bib}

\usepackage[hidelinks]{hyperref}
\hypersetup{
  colorlinks=true,
  linkcolor=accent,
  citecolor=accent,
  urlcolor=accent
}

% ---------- 标题层级 ----------
\ctexset{
  section = {
    format    = \Large\bfseries\heiti,
    beforeskip = 1.2ex plus .2ex,
    afterskip  = .8ex plus .1ex,
  },
  subsection = {
    format    = \large\bfseries\heiti,
    beforeskip = 1.0ex plus .2ex,
    afterskip  = .6ex plus .1ex,
  },
  subsubsection = {
    format    = \normalsize\bfseries\heiti,
    beforeskip = .8ex plus .2ex,
    afterskip  = .4ex plus .1ex,
  },
}

% ---------- 题注 ----------
\captionsetup{font=small,labelsep=quad}
\captionsetup[table]{position=top,skip=4pt}
\captionsetup[figure]{position=bottom,skip=4pt}

% ---------- 列表 ----------
\setlist{nosep,leftmargin=2em}

% ---------- 摘要环境 ----------
\newenvironment{cnabstract}{%
  \begin{center}\large\bfseries\heiti 摘要\end{center}
  \vspace{-6pt}
  \begin{quote}
}{%
  \end{quote}
}

\begin{document}
\linespread{1.35}\selectfont
\pagestyle{plain}

% ============================================================
% 标题区
% ============================================================
\begin{center}
  {\LARGE\bfseries\heiti 论文题目：在此填写标题\par}
  \vspace{12pt}
  {\normalsize 作者姓名\par}
  \vspace{3pt}
  {\small 单位名称，城市 邮编\par}
  \vspace{3pt}
  {\small \today\par}
\end{center}

\vspace{6pt}

% ============================================================
% 摘要
% ============================================================
\begin{cnabstract}
本文针对\ldots\ldots 问题，提出了\ldots\ldots 方法。该方法通过\ldots\ldots 实现\ldots\ldots。
实验结果表明，所提方法在\ldots\ldots 数据集上将\ldots\ldots 指标提升至\ldots\ldots（给出具体数值），
较基线方法提高\ldots\ldots。结论适用于\ldots\ldots 条件。
\end{cnabstract}

\noindent\textbf{关键词：}关键词一；关键词二；关键词三

\vspace{8pt}

% 英文摘要（按需保留或删除）
\begin{center}\large\bfseries Abstract\end{center}
\vspace{-6pt}
\begin{quote}
\small
This paper addresses \ldots\ldots. We propose \ldots\ldots. Experiments show that \ldots\ldots.
\end{quote}
\noindent\textbf{Keywords:} keyword one; keyword two; keyword three

\vspace{10pt}

% ============================================================
% 正文
% ============================================================
\section{引言}
\label{sec:intro}

领域背景与研究意义\cite{ref1}。已有工作主要沿两条路线展开\cite{ref2,ref3}，但尚未解决\ldots\ldots。
本文针对该问题提出\ldots\ldots。

本文的主要贡献如下：
\begin{enumerate}
  \item 提出\ldots\ldots 方法，解决\ldots\ldots 问题；
  \item 在\ldots\ldots 数据集上验证\ldots\ldots；
  \item 给出\ldots\ldots 的分析与适用边界。
\end{enumerate}

\section{方法}
\label{sec:method}

\subsection{问题定义}

设输入为 $x \in \mathbb{R}^{d}$，输出为 $y$，目标是最小化如下损失：
\begin{equation}
  \mathcal{L}(\theta) = \frac{1}{N}\sum_{i=1}^{N}\left(y_i - f(x_i;\theta)\right)^2 ,
  \label{eq:loss}
\end{equation}
其中 $N$ 为样本数，$\theta$ 为模型参数，$f(\cdot;\theta)$ 为预测函数。

\subsection{算法流程}

算法步骤如表~\ref{tab:example} 所示。

% ============================================================
% 表格示例（三线表）
% ============================================================
\begin{table}[htbp]
  \centering
  \caption{三种方法的性能对比}
  \label{tab:example}
  \begin{tabular}{lccc}
    \toprule
    方法 & 准确率（\%） & 召回率（\%） & 耗时（s） \\
    \midrule
    基线方法 & 82.1 & 79.4 & 12.3 \\
    对比方法 & 86.5 & 83.1 & 11.6 \\
    本文方法 & 89.6 & 87.2 & 10.8 \\
    \bottomrule
  \end{tabular}
\end{table}

表~\ref{tab:example} 表明，本文方法在准确率与召回率上均优于基线，且耗时更低。

% ============================================================
% 图片示例（替换为真实图片）
% ============================================================
\begin{figure}[htbp]
  \centering
  % \includegraphics[width=0.8\textwidth]{figs/result.pdf}
  \fbox{\rule{0pt}{4cm}\rule{8cm}{0pt}}
  \caption{示例图：图中应传达一个明确信息}
  \label{fig:example}
\end{figure}

如图~\ref{fig:example} 所示，\ldots\ldots。

\section{实验结果与分析}
\label{sec:exp}

\subsection{实验设置}

数据来源、划分方式、评价指标与运行环境在此说明。

\subsection{主结果}

结合表~\ref{tab:example} 与图~\ref{fig:example} 分析\ldots\ldots。

\subsection{消融与误差分析}

\section{结论}
\label{sec:conclusion}

本文提出\ldots\ldots，在\ldots\ldots 条件下取得\ldots\ldots 的结果。
局限性在于\ldots\ldots，后续可从\ldots\ldots 方向改进。

% ============================================================
% 参考文献（默认：biblatex + biber + GB/T 7714-2015）
% ============================================================
\printbibliography[heading=bibintoc,title={参考文献}]

% 方案二：≤ 10 条的小论文可改用手写条目（无需 biber）
% 删掉上面这行 \printbibliography 与导言区的 biblatex 两行，改用：
% \begin{thebibliography}{99}
% \bibitem{ref1} 作者一, 作者二. 中文文献题名[J]. 期刊名称, 2024, 12(3): 45--52.
% \bibitem{ref2} Author A, Author B. English paper title[J]. Journal Name, 2023, 45(2): 100--115.
% \bibitem{ref3} 作者三. 会议论文题名[C]//会议名称. 城市: 出版社, 2023: 1--8.
% \end{thebibliography}

% ============================================================
% 附录（按需）
% ============================================================
\appendix
\section{补充材料}
\label{app:extra}

此处放置推导、代码说明或补充表格。

\end{document}
